% GENERATED FILE. Edit canonical lessons, then run npm run generate:books.
% canonical-source-hash: fnv1a64:58da3abdc0420fb1
% canonical-lessons: KA-C334-bananti, KA-C334-janana, KA-C334-namakarana, KA-C334-simanta, KA-C334-uttaradhikari

\chapter{A Baby in the Family}
\label{ch:ka-a-baby-in-the-family}
\hlchaptermodality{334}

\begin{chapteropening}
\textbf{By the end of this chapter:} \emph{I can say a new mother, a birth, a naming ceremony, a baby shower and a successor.}

Complete the last lesson of chapter 334: I can say a new mother, a birth, a naming ceremony, a baby shower and a successor.
\end{chapteropening}

\section[\texorpdfstring{\kn{ಬಾಣಂತಿ} (bāṇanti)}{ಬಾಣಂತಿ (bāṇanti)}]{\kn{ಬಾಣಂತಿ} (bāṇanti) — a new mother}

\label{lesson:KA-C334-bananti}



\begin{quote}
Before the new one: say the Kannada for to set up, then the Kannada for an invitation card.
\end{quote}

\subsection*{You'll want to know: \kn{ಬಾಣಂತಿ}}
\textbf{\kn{ಬಾಣಂತಿ}} — \emph{bāṇanti} — \textquotedblleft{}a new mother\textquotedblright{}.

\begin{quote}
\kn{ಬಾಣಂತಿ} \kn{ಮೂರು} \kn{ತಿಂಗಳು} \kn{ಅಮ್ಮನ} \kn{ಮನೆಯಲ್ಲಿ} \kn{ಇದ್ದಳು}.
\end{quote}

\begin{quote}
\emph{bāṇanti mūru tiṅgaḷu ammana maneyalli iddaḷu.}
\end{quote}

\begin{quote}
The new mother stayed three months at her mother's house.
\end{quote}

A woman in the weeks after giving birth, looked after with rest, warm food and the home-made \kn{ಔಷಧ}, medicine, that the elders prepare.

\begin{cousinweb}
\textbf{\kn{ಬಾಣಂತಿ}} (\emph{bāṇanti}) is native Dravidian: DEDR 5347 lists Kannada \emph{bāṇanti}, \textquotedblleft{}a woman recently delivered\textquotedblright{}, beside Malayalam \emph{vaṇṇa}.
\end{cousinweb}

\begin{grammarlens}[title={What you've built}]
One more word for telling what happened.
\end{grammarlens}

\subsection*{Your turn}
\begin{itemize}
\raggedright
  \item \emph{Say it:} \emph{bāṇanti}
  \item \emph{Say it:} \emph{bāṇanti}, once more
  \item \emph{Say it:} \emph{bāṇanti mūru tiṅgaḷu ammana maneyalli iddaḷu.} — The new mother stayed three months at her mother's house.
  \item \emph{Recall:} say the Kannada for to set up, then the Kannada for an invitation card, then say \emph{bāṇanti} again
\end{itemize}

\begin{tcolorbox}[breakable,colback=teal!4,colframe=teal!35!black,title={Before you move on}]
What does \kn{ಬಾಣಂತಿ} mean? (\textquotedblleft{}A new mother\textquotedblright{}.) Say it once more.
\end{tcolorbox}
\section[\texorpdfstring{\kn{ಜನನ} (janana)}{ಜನನ (janana)}]{\kn{ಜನನ} (janana) — a birth}

\label{lesson:KA-C334-janana}



\begin{quote}
Before the new one: say the Kannada for an invitation card, then the Kannada for a new mother.
\end{quote}

\subsection*{You'll want to know: \kn{ಜನನ}}
\textbf{\kn{ಜನನ}} — \emph{janana} — \textquotedblleft{}a birth\textquotedblright{}.

\begin{quote}
\kn{ಮಗುವಿನ} \kn{ಜನನ} \kn{ಆಸ್ಪತ್ರೆಯಲ್ಲಿ} \kn{ಆಯಿತು}.
\end{quote}

\begin{quote}
\emph{maguvina janana āspatreyalli āyitu.}
\end{quote}

\begin{quote}
The baby was born in the hospital.
\end{quote}

A formal word: \kn{ಜನನ} \kn{ಪ್ರಮಾಣಪತ್ರ}, a birth certificate. At home people say \kn{ಹುಟ್ಟು}.

\begin{cousinweb}
\textbf{\kn{ಜನನ}} (\emph{janana}) is borrowed from Sanskrit \emph{janana}, \textquotedblleft{}birth\textquotedblright{} (CDIAL 5099), from the root \emph{jan}, \textquotedblleft{}to be born\textquotedblright{}.
\end{cousinweb}

\begin{grammarlens}[title={What you've built}]
One more word for telling what happened.
\end{grammarlens}

\subsection*{Your turn}
\begin{itemize}
\raggedright
  \item \emph{Say it:} \emph{janana}
  \item \emph{Say it:} \emph{janana}, once more
  \item \emph{Say it:} \emph{maguvina janana āspatreyalli āyitu.} — The baby was born in the hospital.
  \item \emph{Recall:} say the Kannada for an invitation card, then the Kannada for a new mother, then say \emph{janana} again
\end{itemize}

\begin{tcolorbox}[breakable,colback=teal!4,colframe=teal!35!black,title={Before you move on}]
What does \kn{ಜನನ} mean? (\textquotedblleft{}A birth\textquotedblright{}.) Say it once more.
\end{tcolorbox}
\section[\texorpdfstring{\kn{ನಾಮಕರಣ} (nāmakaraṇa)}{ನಾಮಕರಣ (nāmakaraṇa)}]{\kn{ನಾಮಕರಣ} (nāmakaraṇa) — a naming ceremony}

\label{lesson:KA-C334-namakarana}



\begin{quote}
Before the new one: say the Kannada for a new mother, then the Kannada for a birth.
\end{quote}

\subsection*{You'll want to know: \kn{ನಾಮಕರಣ}}
\textbf{\kn{ನಾಮಕರಣ}} — \emph{nāmakaraṇa} — \textquotedblleft{}a naming ceremony\textquotedblright{}.

\begin{quote}
\kn{ಮಗುವಿನ} \kn{ನಾಮಕರಣ} \kn{ಹನ್ನೊಂದನೆಯ} \kn{ದಿನ} \kn{ನಡೆಯಿತು}.
\end{quote}

\begin{quote}
\emph{maguvina nāmakaraṇa hannondaneya dina naḍeyitu.}
\end{quote}

\begin{quote}
The baby's naming ceremony took place on the eleventh day.
\end{quote}

The day a baby gets its name, often with the baby in a cradle.

\begin{cousinweb}
\textbf{\kn{ನಾಮಕರಣ}} (\emph{nāmakaraṇa}) is borrowed from Sanskrit \emph{nāmakaraṇa}, \textquotedblleft{}the ceremony of naming after birth\textquotedblright{} (CDIAL 7066), from \emph{nāman} \textquotedblleft{}a name\textquotedblright{} and \emph{karaṇa} \textquotedblleft{}making\textquotedblright{}.
\end{cousinweb}

\begin{grammarlens}[title={What you've built}]
One more word for telling what happened.
\end{grammarlens}

\subsection*{Your turn}
\begin{itemize}
\raggedright
  \item \emph{Say it:} \emph{nāmakaraṇa}
  \item \emph{Say it:} \emph{nāmakaraṇa}, once more
  \item \emph{Say it:} \emph{maguvina nāmakaraṇa hannondaneya dina naḍeyitu.} — The baby's naming ceremony took place on the eleventh day.
  \item \emph{Recall:} say the Kannada for a new mother, then the Kannada for a birth, then say \emph{nāmakaraṇa} again
\end{itemize}

\begin{tcolorbox}[breakable,colback=teal!4,colframe=teal!35!black,title={Before you move on}]
What does \kn{ನಾಮಕರಣ} mean? (\textquotedblleft{}A naming ceremony\textquotedblright{}.) Say it once more.
\end{tcolorbox}
\section[\texorpdfstring{\kn{ಸೀಮಂತ} (sīmanta)}{ಸೀಮಂತ (sīmanta)}]{\kn{ಸೀಮಂತ} (sīmanta) — a baby shower}

\label{lesson:KA-C334-simanta}



\begin{quote}
Before the new one: say the Kannada for a birth, then the Kannada for a naming ceremony.
\end{quote}

\subsection*{You'll want to know: \kn{ಸೀಮಂತ}}
\textbf{\kn{ಸೀಮಂತ}} — \emph{sīmanta} — \textquotedblleft{}a baby shower\textquotedblright{}.

\begin{quote}
\kn{ಅಕ್ಕನ} \kn{ಸೀಮಂತಕ್ಕೆ} \kn{ತುಂಬಾ} \kn{ಹೆಂಗಸರು} \kn{ಬಂದರು}.
\end{quote}

\begin{quote}
\emph{akkana sīmantakke tumbā heṅgasaru bandaru.}
\end{quote}

\begin{quote}
Many women came to my elder sister's baby shower.
\end{quote}

A ceremony for a woman expecting her first child, with bangles, flowers and sweets.

\begin{cousinweb}
\textbf{\kn{ಸೀಮಂತ}} (\emph{sīmanta}) is borrowed from Sanskrit \emph{sīmanta}, \textquotedblleft{}the parting of the hair\textquotedblright{} (CDIAL 13436): the old rite for a pregnant woman began with parting her hair.
\end{cousinweb}

\begin{grammarlens}[title={What you've built}]
One more word for telling what happened.
\end{grammarlens}

\subsection*{Your turn}
\begin{itemize}
\raggedright
  \item \emph{Say it:} \emph{sīmanta}
  \item \emph{Say it:} \emph{sīmanta}, once more
  \item \emph{Say it:} \emph{akkana sīmantakke tumbā heṅgasaru bandaru.} — Many women came to my elder sister's baby shower.
  \item \emph{Recall:} say the Kannada for a birth, then the Kannada for a naming ceremony, then say \emph{sīmanta} again
\end{itemize}

\begin{tcolorbox}[breakable,colback=teal!4,colframe=teal!35!black,title={Before you move on}]
What does \kn{ಸೀಮಂತ} mean? (\textquotedblleft{}A baby shower\textquotedblright{}.) Say it once more.
\end{tcolorbox}
\section[\texorpdfstring{\kn{ಉತ್ತರಾಧಿಕಾರಿ} (uttarādhikāri)}{ಉತ್ತರಾಧಿಕಾರಿ (uttarādhikāri)}]{\kn{ಉತ್ತರಾಧಿಕಾರಿ} (uttarādhikāri) — a successor}

\label{lesson:KA-C334-uttaradhikari}



\begin{quote}
Before the new one: say the Kannada for a naming ceremony, then the Kannada for a baby shower.
\end{quote}

\subsection*{You'll want to know: \kn{ಉತ್ತರಾಧಿಕಾರಿ}}
\textbf{\kn{ಉತ್ತರಾಧಿಕಾರಿ}} — \emph{uttarādhikāri} — \textquotedblleft{}a successor\textquotedblright{}.

\begin{quote}
\kn{ಆ} \kn{ಮಗು} \kn{ರಾಜನ} \kn{ಉತ್ತರಾಧಿಕಾರಿ} \kn{ಆಯಿತು}.
\end{quote}

\begin{quote}
\emph{ā magu rājana uttarādhikāri āyitu.}
\end{quote}

\begin{quote}
That child became the king's successor.
\end{quote}

Plural: \kn{ಉತ್ತರಾಧಿಕಾರಿಗಳು}, \emph{uttarādhikārigaḷu}.

The one who takes over a throne, a post or property after someone.

\begin{cousinweb}
\textbf{\kn{ಉತ್ತರಾಧಿಕಾರಿ}} (\emph{uttarādhikāri}) is borrowed from Sanskrit \emph{uttarādhikārin}, \textquotedblleft{}an heir or claimant after the death of the owner\textquotedblright{} (Monier-Williams).
\end{cousinweb}

\begin{grammarlens}[title={What you've built}]
One more word for telling what happened.
\end{grammarlens}

\subsection*{Your turn}
\begin{itemize}
\raggedright
  \item \emph{Say it:} \emph{uttarādhikāri}
  \item \emph{Say it:} \emph{uttarādhikāri}, once more
  \item \emph{Say it:} \emph{ā magu rājana uttarādhikāri āyitu.} — That child became the king's successor.
  \item \emph{Recall:} say the Kannada for a naming ceremony, then the Kannada for a baby shower, then say \emph{uttarādhikāri} again
\end{itemize}

\begin{tcolorbox}[breakable,colback=teal!4,colframe=teal!35!black,title={Before you move on}]
What does \kn{ಉತ್ತರಾಧಿಕಾರಿ} mean? (\textquotedblleft{}A successor\textquotedblright{}.) Say it once more.
\end{tcolorbox}
